% ============================================
% INTRODUCTION
% ============================================
% 👤 PERSON 2: Background & Threat Modeling
% 
% Instructions:
% - Start with ransomware impact/motivation (statistics, real incidents)
% - Identify the gap: why existing solutions fail
% - Introduce ADRDS as the solution
% - Include the KILLER DEMO PARAGRAPH here or in Demo section
% - List paper contributions clearly (3-4 bullets)
% - Length: 0.8-1 page
% ============================================

\section{Introduction}

[PERSON 2: Set the stage - why ransomware matters and why ADRDS matters]

Ransomware attacks have escalated dramatically in recent years, causing billions 
of dollars in damages and disrupting critical infrastructure worldwide. Traditional 
signature-based detection methods struggle to keep pace with rapidly evolving 
ransomware variants, creating an urgent need for adaptive, behavior-driven approaches.

\subsection{Motivation}

[Ransomware remains one of the most pervasive and damaging cyber threats. According to the 2025 Verizon Data Breach Investigations Report, ransomware accounted for over 25% of all malware-related incidents, with global damages exceeding $30 billion in 2024 alone. High-profile attacks such as the 2024 MedStar Health breach and the Colonial Pipeline incident have demonstrated the ability of ransomware to disrupt critical services, compromise sensitive data, and inflict severe financial losses. Attackers increasingly leverage double extortion tactics, demanding payment not only for decryption but also to prevent data leaks. Despite increased investment in cybersecurity, organizations continue to face challenges in detecting and responding to rapidly evolving ransomware variants, which often bypass traditional signature-based defenses.]

\subsection{Research Gap}

Current ransomware detection solutions predominantly rely on static signatures or heuristic rules, which are easily evaded by novel or obfuscated ransomware strains. While machine learning-based approaches have shown promise, most are evaluated in offline settings and lack transparency or interactivity for end-users. There is a critical need for systems that not only detect ransomware dynamically but also provide real-time feedback, visualization, and user-driven exploration. Such interactive platforms are essential for security education, research, and rapid response, yet remain largely absent in existing literature and commercial offerings.

\subsection{Proposed Solution}

We present ADRDS (Adaptive Dynamic Ransomware Detection System), an interactive 
demonstration platform that combines behavioral taxonomy, machine learning 
classification, and real-time visualization. \textbf{During the live demo, 
attendees can interactively control ransomware stages, observe real-time behavioral 
visualization, and see on-the-fly ML classification outcomes with adjustable thresholds.}

\subsection{Contributions}

This paper makes the following contributions:
\begin{itemize}
    \item A comprehensive behavioral taxonomy for modeling ransomware attack stages
    \item An interactive demonstration system integrating simulation, detection, and visualization
    \item Experimental evaluation on [dataset] achieving [X\%] detection accuracy
    \item An open platform for security education and research
\end{itemize}

The remainder of this paper is organized as follows: Section II reviews related work, 
Section III defines our threat model, Sections IV-V describe the system architecture 
and taxonomy, Section VI details implementation, Section VII presents the demo scenarios, 
Sections VIII-X provide experimental evaluation, and Section XI concludes.
